% ============================================================
\chapter{Symmetric Spaces}
\label{ch:symmetric}
% ============================================================

\'Elie Cartan, in the 1920s--1930s, undertook one of the most ambitious classifications in geometry: that of Riemannian spaces whose curvature is ``the same at every point and in every direction'' in the sense that the geodesic symmetry at each point is a global isometry. These \emph{symmetric spaces} encompass the spaces of constant curvature (spheres, hyperbolic spaces), projective spaces, Grassmannians, and many others. Cartan showed that their classification reduces to that of semi-simple Lie algebras, establishing a deep bridge between geometry and algebra. Today, symmetric spaces appear in theoretical physics, number theory, and machine learning (optimization on Stiefel and Grassmann manifolds).

\section{Definition and first examples}

\begin{definition}[Riemannian symmetric space]
A complete Riemannian manifold $(M, g)$ is a \emph{Riemannian symmetric space}
if for every point $p \in M$, there exists an involutive isometry
$s_p : M \to M$ such that:
\begin{enumerate}
    \item $s_p(p) = p$ (fixed point).
    \item $(ds_p)_p = -\operatorname{Id}_{T_pM}$ (the differential is $-\operatorname{Id}$).
\end{enumerate}
The isometry $s_p$ is called the \emph{geodesic symmetry} at $p$.
\end{definition}

\begin{remark}
The condition $(ds_p)_p = -\operatorname{Id}$ implies that $s_p$ reverses all
geodesics through $p$: $s_p(\gamma(t)) = \gamma(-t)$.
\end{remark}

\begin{example}[Model spaces]
\begin{itemize}
    \item $\R^n$: $s_p(x) = 2p - x$ (central symmetry).
    \item $S^n$: $s_p$ is the reflection through the subspace passing through
          $p$ and the center.
    \item $\mathbb{H}^n$: in the disk model, $s_0$ is the inversion $x \mapsto -x$.
\end{itemize}
\end{example}

\begin{example}[Projective spaces]
$\R P^n$, $\C P^n$, $\mathbb{H}P^n$ (quaternionic projective space) and the
Cayley plane $\mathbb{O}P^2$ are compact symmetric spaces of rank $1$.
\end{example}

\begin{example}[Compact Lie groups]
Every compact Lie group $G$ with a bi-invariant metric is a symmetric space.
The symmetry at the identity is $s_e(g) = g^{-1}$.
\end{example}

\section{Symmetric pair structure}

\begin{definition}[Symmetric pair]
A \emph{Riemannian symmetric pair} is a couple $(G, K)$ where:
\begin{itemize}
    \item $G$ is a connected Lie group.
    \item $K$ is a compact subgroup of $G$.
    \item There exists an involutive automorphism $\sigma : G \to G$ such that
          $(G^\sigma)_0 \subset K \subset G^\sigma$, where $G^\sigma$ is the
          fixed-point subgroup of $\sigma$.
\end{itemize}
\end{definition}

\begin{theorem}[Fundamental correspondence]
There is a bijective correspondence between:
\begin{enumerate}
    \item Simply connected symmetric spaces $(M, g)$.
    \item Effective symmetric pairs $(G, K, \sigma)$ with $G$ simply connected.
\end{enumerate}
The symmetric space is $M = G/K$ and the metric is induced by an
$\operatorname{Ad}(K)$-invariant inner product on $\mathfrak{p}$.
\end{theorem}

\section{Cartan decomposition}

\begin{definition}[Cartan decomposition]
Let $\sigma : G \to G$ be the involution and $\sigma_* : \mathfrak{g} \to \mathfrak{g}$
its differential. The Lie algebra decomposes into eigenspaces:
\[
    \mathfrak{g} = \mathfrak{k} \oplus \mathfrak{p},
\]
where $\mathfrak{k} = \{X \in \mathfrak{g} : \sigma_*(X) = X\}$ ($+1$ eigenspace)
and $\mathfrak{p} = \{X \in \mathfrak{g} : \sigma_*(X) = -X\}$ ($-1$ eigenspace).
\end{definition}

\begin{proposition}[Bracket relations]
The Lie brackets satisfy:
\[
    [\mathfrak{k}, \mathfrak{k}] \subset \mathfrak{k}, \quad
    [\mathfrak{k}, \mathfrak{p}] \subset \mathfrak{p}, \quad
    [\mathfrak{p}, \mathfrak{p}] \subset \mathfrak{k}.
\]
\end{proposition}

\begin{proof}
If $X, Y \in \mathfrak{k}$, then $\sigma_*[X,Y] = [\sigma_*X, \sigma_*Y] = [X,Y]$,
so $[X,Y] \in \mathfrak{k}$. If $X \in \mathfrak{k}$, $Y \in \mathfrak{p}$,
then $\sigma_*[X,Y] = [X,-Y] = -[X,Y]$, so $[X,Y] \in \mathfrak{p}$.
If $X, Y \in \mathfrak{p}$, then $\sigma_*[X,Y] = [-X,-Y] = [X,Y]$,
so $[X,Y] \in \mathfrak{k}$.
\end{proof}

The tangent space to $M = G/K$ at the base point $o = eK$ is identified with
$\mathfrak{p}$, and the Riemannian metric is an $\operatorname{Ad}(K)$-invariant
inner product on $\mathfrak{p}$.

\section{Curvature of symmetric spaces}

\begin{theorem}[Curvature formula]
On the symmetric space $M = G/K$, the Levi-Civita connection satisfies, for
$X, Y \in \mathfrak{p} \cong T_oM$:
\[
    \nabla_X Y = \frac{1}{2}[X, Y]_\mathfrak{p} = 0,
\]
since $[X, Y] \in \mathfrak{k}$ and the projection onto $\mathfrak{p}$ vanishes.
The curvature tensor is:
\[
    R(X, Y)Z = -[[X, Y], Z], \quad X, Y, Z \in \mathfrak{p}.
\]
\end{theorem}

\begin{corollary}
The curvature of a symmetric space is entirely determined by the Lie brackets.
In particular, $\nabla R = 0$ (the curvature is parallel).
\end{corollary}

\begin{theorem}[Characterization by $\nabla R = 0$]
A complete simply connected Riemannian manifold $(M, g)$ is a symmetric space
if and only if $\nabla R = 0$.
\end{theorem}

\begin{proposition}[Sectional curvature]
For orthonormal $X, Y \in \mathfrak{p}$:
\[
    K(X, Y) = \inner{[[X,Y], X]}{Y} = \norm{[X, Y]}^2 \geq 0
\]
if the metric comes from the negative Killing form (compact type).
For noncompact type, $K(X,Y) = -\norm{[X,Y]}^2 \leq 0$.
\end{proposition}

\section{Classification of symmetric spaces}

\begin{theorem}[de Rham decomposition]
Every simply connected symmetric space decomposes uniquely as a product:
\[
    M = M_0 \times M_1 \times \cdots \times M_k,
\]
where $M_0 = \R^p$ is the Euclidean factor and the $M_i$ ($i \geq 1$) are
irreducible symmetric spaces.
\end{theorem}

\begin{definition}[Types of symmetric spaces]
Irreducible symmetric spaces are classified into:
\begin{itemize}
    \item \textbf{Compact type:} $K \geq 0$, $G$ compact. Example: $S^n$, $\C P^n$.
    \item \textbf{Noncompact type:} $K \leq 0$, $G$ noncompact.
          Example: $\mathbb{H}^n$, $SL(n,\R)/SO(n)$.
    \item \textbf{Duality:} to every compact space corresponds a ``dual'' noncompact
          space.
\end{itemize}
\end{definition}

\begin{keyformulas}[title=\textbf{Examples of irreducible symmetric spaces}]
\begin{center}
\renewcommand{\arraystretch}{1.3}
\begin{tabular}{lccl}
\hline
\textbf{Space} & $G/K$ & \textbf{Type} & $\dim$ \\
\hline
Sphere $S^n$ & $SO(n+1)/SO(n)$ & compact & $n$ \\
$\C P^n$ & $SU(n+1)/S(U(1) \times U(n))$ & compact & $2n$ \\
Grassmannian & $SO(p+q)/S(O(p) \times O(q))$ & compact & $pq$ \\
$\mathbb{H}^n$ & $SO_0(n,1)/SO(n)$ & noncompact & $n$ \\
$SL(n,\R)/SO(n)$ & & noncompact & $\frac{n(n+1)}{2}-1$ \\
\hline
\end{tabular}
\end{center}
\end{keyformulas}

\section{Rank and flat subspaces}

\begin{definition}[Rank]
The \emph{rank} of a symmetric space is the maximal dimension of an abelian
subspace $\mathfrak{a} \subset \mathfrak{p}$ (i.e.\ $[\mathfrak{a}, \mathfrak{a}] = 0$).
Geometrically, it is the dimension of the largest totally geodesic flat subspace.
\end{definition}

\begin{example}
$S^n$ and $\mathbb{H}^n$ have rank $1$. $SL(n,\R)/SO(n)$ has rank $n-1$.
The Grassmannian $G_{p,q}$ has rank $\min(p,q)$.
\end{example}

\section{Geodesics and group structure}

\begin{proposition}[Geodesics]
The geodesics of $M = G/K$ through $o = eK$ are the curves
$t \mapsto \exp(tX) \cdot o$ for $X \in \mathfrak{p}$.
\end{proposition}

\begin{proposition}[Parallel transport]
Parallel transport along a geodesic $\gamma(t) = \exp(tX) \cdot o$ is given
by the action of $\exp(tX)$ on $T_oM \cong \mathfrak{p}$.
\end{proposition}

\begin{theorem}[Holonomy]
The restricted holonomy group of an irreducible symmetric space is
$\operatorname{Hol}^0(M) = K$ (the connected component). Berger's classification
of holonomy groups of irreducible non-symmetric manifolds complements the
symmetric space classification.
\end{theorem}

\section{Rank-one symmetric spaces}

The compact simply connected rank-one symmetric spaces are:
\begin{enumerate}
    \item $S^n = SO(n+1)/SO(n)$, curvature $K = 1$.
    \item $\C P^n = SU(n+1)/S(U(1) \times U(n))$, curvature $\frac{1}{4} \leq K \leq 1$.
    \item $\mathbb{H}P^n = Sp(n+1)/(Sp(1) \times Sp(n))$, curvature $\frac{1}{4} \leq K \leq 1$.
    \item $\mathbb{O}P^2 = F_4/\operatorname{Spin}(9)$, curvature $\frac{1}{4} \leq K \leq 1$.
\end{enumerate}

\begin{remark}
For $\C P^n$, the sectional curvature varies between $1/4$ and $1$. The minimum
$1/4$ is achieved on totally real planes, and the maximum $1$ on holomorphic planes.
\end{remark}

\section{Exercises}

\begin{exercise}
Show that $S^n$ is a symmetric space by explicitly exhibiting the symmetry $s_p$
at every point $p$.
\end{exercise}

\begin{exercise}
Compute the Cartan decomposition $\mathfrak{g} = \mathfrak{k} \oplus \mathfrak{p}$
for $S^2 = SO(3)/SO(2)$ and verify the bracket relations.
\end{exercise}

\begin{exercise}
Show that for the Lie group $G$ with bi-invariant metric viewed as the symmetric
space $(G \times G)/\operatorname{diag}(G)$, the sectional curvature satisfies
$K(X,Y) = \frac{1}{4}\norm{[X,Y]}^2$.
\end{exercise}

\begin{exercise}
Show that $\C P^n$ is a symmetric space and compute its sectional curvature.
Verify that $1/4 \leq K \leq 1$.
\end{exercise}

\begin{exercise}
Show that $SL(2,\R)/SO(2) \cong \mathbb{H}^2$ by exhibiting an explicit isometry.
\end{exercise}

\begin{exercise}
Determine the rank of the Grassmannian $G_{2,3}(\R) = SO(5)/(SO(2) \times SO(3))$.
\end{exercise}
